% F - input tensor shape/orientation (Chapter 2). Clarifies that each trial is a
% C=6 (channels, IMU axes) by T=30 (timesteps) array, and that the network reads
% ONE column (one timestep) at a time into a LIF neuron.
\documentclass[border=2pt]{standalone}
\usepackage{tikz}
% =====================================================================
%  Shared TikZ style for the BISCCITS SNN workshop figures.
%  Every figure file uses:
%      \documentclass[border=2pt]{standalone}
%      \usepackage{tikz}
%      % =====================================================================
%  Shared TikZ style for the BISCCITS SNN workshop figures.
%  Every figure file uses:
%      \documentclass[border=2pt]{standalone}
%      \usepackage{tikz}
%      % =====================================================================
%  Shared TikZ style for the BISCCITS SNN workshop figures.
%  Every figure file uses:
%      \documentclass[border=2pt]{standalone}
%      \usepackage{tikz}
%      \input{style.tex}
%  so that all figures share one visual language (see SKILL: tikz-figure).
% =====================================================================
\usepackage{amsmath}
\usepackage{amssymb}   % \mathbb, etc.
\usetikzlibrary{arrows.meta, positioning, calc, backgrounds, fit,
                decorations.pathmorphing, shapes.misc}

% ----- palette: matches the matplotlib tab: colours used in the notebook -----
\definecolor{tabblue}{HTML}{1F77B4}    % input / current
\definecolor{taborange}{HTML}{FF7F0E}  % accent / secondary
\definecolor{tabgreen}{HTML}{2CA02C}   % SNN / "good"
\definecolor{tabred}{HTML}{D62728}     % threshold / reset
\definecolor{tabpurple}{HTML}{9467BD}  % membrane potential V
\definecolor{tabgray}{HTML}{7F7F7F}    % leak / inactive / captions
\definecolor{inkblack}{HTML}{222222}   % spikes / outlines / text

% ----- typography -----
\tikzset{every node/.style={font=\small}}

% ----- canonical element styles -----
\tikzset{
  % the LIF neuron glyph -- reuse this EVERYWHERE a neuron appears
  lifnode/.style={circle, draw=inkblack, line width=0.7pt,
                  minimum size=11mm, fill=tabpurple!12, inner sep=1pt},
  % a smaller neuron for dense multi-neuron layouts
  lifdot/.style={circle, draw=inkblack, line width=0.6pt,
                 minimum size=7mm, fill=tabpurple!12, inner sep=0pt},
  iolabel/.style={font=\footnotesize, inner sep=1.5pt},
  current/.style={-{Stealth[length=2.2mm]}, tabblue,  line width=0.9pt}, % input current
  synapse/.style={-{Stealth[length=2.2mm]}, inkblack, line width=0.8pt}, % forward connection
  spike/.style  ={-{Stealth[length=2.4mm]}, inkblack, line width=1.0pt}, % spike output
  leak/.style   ={-{Stealth[length=1.8mm]}, tabgray,  line width=0.8pt, dashed}, % leak / recurrence
  inactive/.style={tabgray!55, line width=0.6pt},      % "skipped" / silent path
  thresholdline/.style={tabred, dashed, line width=0.8pt},
  block/.style={draw=inkblack, rounded corners=2pt, fill=white,
                minimum height=8mm, minimum width=16mm, inner sep=3pt},
  layerbox/.style={draw=tabgray, rounded corners=2pt, dashed, inner sep=3mm},
  fignote/.style={font=\footnotesize\itshape, text=tabgray},
}

% ----- canonical symbols (write them identically in every figure) -----
\newcommand{\Vmem}{V}
\newcommand{\Vthr}{V_{\mathrm{thr}}}
\newcommand{\Vreset}{V_{\mathrm{reset}}}
\newcommand{\Iin}{I}
\newcommand{\Sout}{S}
\newcommand{\bbeta}{\beta}
\newcommand{\Wlayer}{W}
\newcommand{\slope}{k}

%  so that all figures share one visual language (see SKILL: tikz-figure).
% =====================================================================
\usepackage{amsmath}
\usepackage{amssymb}   % \mathbb, etc.
\usetikzlibrary{arrows.meta, positioning, calc, backgrounds, fit,
                decorations.pathmorphing, shapes.misc}

% ----- palette: matches the matplotlib tab: colours used in the notebook -----
\definecolor{tabblue}{HTML}{1F77B4}    % input / current
\definecolor{taborange}{HTML}{FF7F0E}  % accent / secondary
\definecolor{tabgreen}{HTML}{2CA02C}   % SNN / "good"
\definecolor{tabred}{HTML}{D62728}     % threshold / reset
\definecolor{tabpurple}{HTML}{9467BD}  % membrane potential V
\definecolor{tabgray}{HTML}{7F7F7F}    % leak / inactive / captions
\definecolor{inkblack}{HTML}{222222}   % spikes / outlines / text

% ----- typography -----
\tikzset{every node/.style={font=\small}}

% ----- canonical element styles -----
\tikzset{
  % the LIF neuron glyph -- reuse this EVERYWHERE a neuron appears
  lifnode/.style={circle, draw=inkblack, line width=0.7pt,
                  minimum size=11mm, fill=tabpurple!12, inner sep=1pt},
  % a smaller neuron for dense multi-neuron layouts
  lifdot/.style={circle, draw=inkblack, line width=0.6pt,
                 minimum size=7mm, fill=tabpurple!12, inner sep=0pt},
  iolabel/.style={font=\footnotesize, inner sep=1.5pt},
  current/.style={-{Stealth[length=2.2mm]}, tabblue,  line width=0.9pt}, % input current
  synapse/.style={-{Stealth[length=2.2mm]}, inkblack, line width=0.8pt}, % forward connection
  spike/.style  ={-{Stealth[length=2.4mm]}, inkblack, line width=1.0pt}, % spike output
  leak/.style   ={-{Stealth[length=1.8mm]}, tabgray,  line width=0.8pt, dashed}, % leak / recurrence
  inactive/.style={tabgray!55, line width=0.6pt},      % "skipped" / silent path
  thresholdline/.style={tabred, dashed, line width=0.8pt},
  block/.style={draw=inkblack, rounded corners=2pt, fill=white,
                minimum height=8mm, minimum width=16mm, inner sep=3pt},
  layerbox/.style={draw=tabgray, rounded corners=2pt, dashed, inner sep=3mm},
  fignote/.style={font=\footnotesize\itshape, text=tabgray},
}

% ----- canonical symbols (write them identically in every figure) -----
\newcommand{\Vmem}{V}
\newcommand{\Vthr}{V_{\mathrm{thr}}}
\newcommand{\Vreset}{V_{\mathrm{reset}}}
\newcommand{\Iin}{I}
\newcommand{\Sout}{S}
\newcommand{\bbeta}{\beta}
\newcommand{\Wlayer}{W}
\newcommand{\slope}{k}

%  so that all figures share one visual language (see SKILL: tikz-figure).
% =====================================================================
\usepackage{amsmath}
\usepackage{amssymb}   % \mathbb, etc.
\usetikzlibrary{arrows.meta, positioning, calc, backgrounds, fit,
                decorations.pathmorphing, shapes.misc}

% ----- palette: matches the matplotlib tab: colours used in the notebook -----
\definecolor{tabblue}{HTML}{1F77B4}    % input / current
\definecolor{taborange}{HTML}{FF7F0E}  % accent / secondary
\definecolor{tabgreen}{HTML}{2CA02C}   % SNN / "good"
\definecolor{tabred}{HTML}{D62728}     % threshold / reset
\definecolor{tabpurple}{HTML}{9467BD}  % membrane potential V
\definecolor{tabgray}{HTML}{7F7F7F}    % leak / inactive / captions
\definecolor{inkblack}{HTML}{222222}   % spikes / outlines / text

% ----- typography -----
\tikzset{every node/.style={font=\small}}

% ----- canonical element styles -----
\tikzset{
  % the LIF neuron glyph -- reuse this EVERYWHERE a neuron appears
  lifnode/.style={circle, draw=inkblack, line width=0.7pt,
                  minimum size=11mm, fill=tabpurple!12, inner sep=1pt},
  % a smaller neuron for dense multi-neuron layouts
  lifdot/.style={circle, draw=inkblack, line width=0.6pt,
                 minimum size=7mm, fill=tabpurple!12, inner sep=0pt},
  iolabel/.style={font=\footnotesize, inner sep=1.5pt},
  current/.style={-{Stealth[length=2.2mm]}, tabblue,  line width=0.9pt}, % input current
  synapse/.style={-{Stealth[length=2.2mm]}, inkblack, line width=0.8pt}, % forward connection
  spike/.style  ={-{Stealth[length=2.4mm]}, inkblack, line width=1.0pt}, % spike output
  leak/.style   ={-{Stealth[length=1.8mm]}, tabgray,  line width=0.8pt, dashed}, % leak / recurrence
  inactive/.style={tabgray!55, line width=0.6pt},      % "skipped" / silent path
  thresholdline/.style={tabred, dashed, line width=0.8pt},
  block/.style={draw=inkblack, rounded corners=2pt, fill=white,
                minimum height=8mm, minimum width=16mm, inner sep=3pt},
  layerbox/.style={draw=tabgray, rounded corners=2pt, dashed, inner sep=3mm},
  fignote/.style={font=\footnotesize\itshape, text=tabgray},
}

% ----- canonical symbols (write them identically in every figure) -----
\newcommand{\Vmem}{V}
\newcommand{\Vthr}{V_{\mathrm{thr}}}
\newcommand{\Vreset}{V_{\mathrm{reset}}}
\newcommand{\Iin}{I}
\newcommand{\Sout}{S}
\newcommand{\bbeta}{\beta}
\newcommand{\Wlayer}{W}
\newcommand{\slope}{k}


\begin{document}
\begin{tikzpicture}[x=1cm, y=1cm]

  % grid geometry
  \def\nrows{6}          % channels C
  \def\ncols{8}          % shown timesteps (trailing ... => 30)
  \def\cell{0.46}        % cell side
  \def\gx{0}             % grid left edge (x)
  \def\gy{0}             % grid bottom edge (y)
  \pgfmathsetmacro{\gw}{\ncols*\cell}   % grid width
  \pgfmathsetmacro{\gh}{\nrows*\cell}   % grid height
  \def\hilite{5}         % index (0-based) of the highlighted column

  % =================================================================
  % smartwatch IMU source block (left)
  % =================================================================
  \node[block, minimum width=18mm, minimum height=12mm, align=center]
        (imu) at (\gx-3.7, \gy+\gh/2) {smartwatch\\ IMU};
  % arrow enters the grid at its TOP-LEFT corner, clear of the row labels below
  \draw[current] (imu.east) |- (\gx-0.05, \gy+\gh-\cell/2)
        node[pos=0.28, above, iolabel, text=tabblue] {records};

  % =================================================================
  % the C x T grid (matrix of small cells)
  % rows are channels (top row = a_x), columns are timesteps
  % =================================================================
  \foreach \r in {0,...,5}{
    \foreach \c in {0,...,7}{
      % highlight one full column to show "one timestep"
      \ifnum\c=\hilite
        \draw[draw=tabblue, line width=0.7pt, fill=tabblue!30]
          (\gx+\c*\cell, \gy+\r*\cell)
          rectangle (\gx+\c*\cell+\cell, \gy+\r*\cell+\cell);
      \else
        \draw[draw=tabgray, line width=0.4pt, fill=tabblue!8]
          (\gx+\c*\cell, \gy+\r*\cell)
          rectangle (\gx+\c*\cell+\cell, \gy+\r*\cell+\cell);
      \fi
    }
  }
  % trailing dots => up to T=30
  \node[iolabel, text=tabgray] at (\gx+\gw+0.35, \gy+\gh/2) {$\cdots$};

  % ---- row labels (the 6 IMU channels), top -> bottom : ax,ay,az,gx,gy,gz ----
  % grid row 5 (top) = a_x ... grid row 0 (bottom) = g_z
  \foreach \r/\lab in {5/{$a_x$}, 4/{$a_y$}, 3/{$a_z$},
                       2/{$g_x$}, 1/{$g_y$}, 0/{$g_z$}}{
    \node[iolabel, anchor=east] at (\gx-0.12, \gy+\r*\cell+\cell/2) {\lab};
  }
  % input-dimension bracket on the left (C is reserved for #classes, so use x in R^6)
  \node[iolabel, text=tabblue, rotate=90, anchor=south]
        at (\gx-1.55, \gy+\gh/2) {input $x\in\mathbb{R}^{6}$};

  % ---- time axis under the columns ----
  \draw[-{Stealth[length=2.2mm]}, inkblack, line width=0.7pt]
        (\gx, \gy-0.32) -- (\gx+\gw+0.55, \gy-0.32);
  \node[iolabel, anchor=west] at (\gx+\gw+0.6, \gy-0.32)
        {time $\rightarrow$ ($T{=}30$)};

  % =================================================================
  % one highlighted column -> a single LIF neuron
  % =================================================================
  \node[lifnode] (n) at (\gx+\gw+3.4, \gy+\gh/2) {$\Vmem[t]$};

  % arrow from the highlighted column into the neuron
  \pgfmathsetmacro{\hx}{\gx+\hilite*\cell+\cell}   % right edge of hilite col
  \draw[current] (\hx, \gy+\gh/2) .. controls
        (\hx+1.3,\gy+\gh/2) and (n.west |- {0,\gy+\gh/2}) .. (n.west)
        node[midway, above, iolabel, text=tabblue, align=center]
        {fed one timestep\\ (column) at a time};

  % spike out of the neuron (small, to show it then runs as usual)
  \draw[spike] (n.east) -- ($(n.east)+(1.0,0)$)
        node[midway, above, iolabel] {$s[t]$};

  % =================================================================
  % fignote
  % =================================================================
  \node[fignote, align=center, text width=12cm]
        at (\gx+\gw/2+0.6, \gy-1.25)
       {each trial is a $6\times 30$ array (input $x$: $6$ channels $\times$
        $T{=}30$ timesteps); the network reads one column per timestep.};

\end{tikzpicture}
\end{document}
